% !TeX encoding = UTF-8
% !TeX spellcheck = pt_BR
% !TeX root = 20201-CODDISC-JOAO-FERREIRA.tex
% ------------------------------------------------------------------------------

% ------------------------------------------------------------------------------
% Template para aulas da UNINASSAU
% Autor: João Ferreira da Silva Júnior
% Versão: 0.1 - 2020-03-27
% ------------------------------------------------------------------------------



% ------------------------------------------------------------------------------
\documentclass[
  % ----------------------------------------------------------------------------
  11pt,
  % ----------------------------------------------------------------------------
  % aspectratio=43,
  aspectratio=169,
  % ----------------------------------------------------------------------------
  % Opções do pacote babel
  english,                    % idioma adicional para hifenização
  french,                     % idioma adicional para hifenização
  spanish,                    % idioma adicional para hifenização
  brazil                      % o último idioma é o principal do documento
  ]{beamer}
% ------------------------------------------------------------------------------





% ------------------------------------------------------------------------------
% ------------------------------------------------------------------------------
% ------------------------------------------------------------------------------
% ------------------------------------------------------------------------------
% ------------------------------------------------------------------------------
% Configurações visuais e de tema

% Extend beamer and a0poster for custom sized posters
%\usepackage[%
%  orientation=landscape,
%  size=custom,
%  width=16,
%  height=9,
%  scale=0.5
%  ]{beamerposter}

% Configura o avanço do slide para passos a cada objeto ou parágrafo do frame
\usepackage{beamerthemesplit}

% This command allows you to specify in a quite general way how a covered item should be rendered
\setbeamercovered{transparent}

% Navigation symbols are small icons that are shown on every slide by default
\setbeamertemplate{navigation symbols}{}

% Configura o tema e a cor
\usetheme{AnnArbor}
\usecolortheme{default}
\usefonttheme{default}
% ------------------------------------------------------------------------------
% Definição de cores
\definecolor{black}{rgb}{0.0, 0.0, 0.0}
\definecolor{white}{rgb}{255,255,255}
\definecolor{red}{rgb}{255,0,0}
\definecolor{green}{rgb}{0,255,0}
\definecolor{blue}{rgb}{0,0,255}
% ------------------------------------------------------------------------------
% Define a aparência dos blocks
\setbeamertemplate{blocks}[default]
% ------------------------------------------------------------------------------



% ------------------------------------------------------------------------------
% ------------------------------------------------------------------------------
% ------------------------------------------------------------------------------
% ------------------------------------------------------------------------------
% ------------------------------------------------------------------------------
% Pacotes principais

% accept different input encodings (conversão automática dos acentos)
\usepackage[utf8]{inputenc}

% standard package for selecting font encodings
\usepackage[T1]{fontenc}
% ------------------------------------------------------------------------------



% ------------------------------------------------------------------------------
% ------------------------------------------------------------------------------
% ------------------------------------------------------------------------------
% ------------------------------------------------------------------------------
% ------------------------------------------------------------------------------
% Pacotes extras

% the package provides several macros to adjust boxed content
\usepackage{adjustbox}

% an extended set of fonts for use in mathematics
\usepackage{amsfonts}

% extension package for LaTeX that provides various features to facilitate writing math formulas and to improve the typographical quality of their output
\usepackage{amsmath}

% math symbols defined by LaTeX package
\usepackage{amssymb}

% Create PDF and SVG animations from graphics files and inline graphics
% \usepackage{animate}

% draw dash-lines in array/tabular environments
% \usepackage{arydshln}

% manages culturally-determined typographical (and other) rules, and hyphenation patterns for a wide range of languages
\usepackage[english,portuguese]{babel}
  \newcommand{\ingles}[1]{\foreignlanguage{english}{\emph{#1}}}
  \newcommand{\frances}[1]{\foreignlanguage{french}{\emph{#1}}}
  % \newcommand{\grego}[1]{\foreignlanguage{greek}{\emph{#1}}}
  \newcommand{\grego}[1]{\emph{#1}}
  \newcommand{\espanhol}[1]{\foreignlanguage{spanish}{\emph{#1}}}
  \newcommand{\portugues}[1]{\foreignlanguage{brazil}{\emph{#1}}}

% a symbol (dingbat) font and LaTeX macros for its use
% \usepackage{bbding}

% enhances the quality of tables in LATEX, providing extra commands as well as behind-the-scenes optimisation
\usepackage{booktabs}

% adding colors to your text is supported by the color package
\usepackage{color}

% gives you Courier scaled at its normal size as the default Typewriter Font
% \usepackage{courier}

% this package provides user control over the layout of the three basic list environments: enumerate, itemize and description
% http://linorg.usp.br/CTAN/macros/latex/contrib/enumitem/enumitem.pdf
% \usepackage{enumitem}

% enhanced support for graphics
\usepackage{graphicx}

% is used to handle cross-referencing commands in LaTeX to produce hypertext links in the document
\usepackage{hyperref}

% this package gives you easy access to the Lorem Ipsum dummy text
\usepackage{lipsum}

% enables the user to typeset programs (programming code) within LATEX
\usepackage{listings}
  \lstset{postbreak=\raisebox{0ex}[0ex][0ex]{\ensuremath{\color{gray}\hookrightarrow\space}}}
  % Aqui tem um exemplo: https://github.com/abntex/abntex2/wiki/HowToFormatarCodigoFonte
  % https://tex.stackexchange.com/questions/64839/how-to-change-listing-caption

  % Listing -> Algorithm
  \renewcommand{\lstlistingname}{Algoritmo}

  % List of Listings -> List of Algorithms
  \renewcommand{\lstlistlistingname}{Lista de algoritmos}

% the Latin Modern family of fonts
\usepackage{lmodern}

% Multimedia inclusion package with Adobe Reader-9/X compatibility
% http://repositorios.cpai.unb.br/ctan/macros/latex/contrib/media9/doc/media9.pdf
\usepackage{media9}

% subliminal refinements towards typographical perfection
\usepackage{microtype}

% Multimedia inclusion package
% Substituído pelo media9
%\usepackage{movie15}
% \newcommand{\pdfmovie}[4]{\href{run:#1}{\framebox{\parbox[c][#3][c]{#2}{\center #4}}}}

% multicol defines a multicols environment which typesets text in multiple columns (up to a maximum of 10)
\usepackage{multicol}

% create tabular cells spanning multiple rows
\usepackage{multirow}

% provides some new list environments. Itemized and enumerated lists can be typesetted within paragraphs
\usepackage{paralist}

% draws high--quality function plots in normal or logarithmic scaling with a user-friendly interface directly in TeX
\usepackage{pgfplots}

% This package reads tab-separated numerical tables from input and generates code for pretty-printed LATEX-tabulars
\usepackage{pgfplotstable}
  \pgfplotsset{compat=1.14}

% provides the means to draw pie (and variant) charts, using PGF/TikZ
\usepackage{pgf-pie}

% this package simplifies the inclusion of external multi-page PDF documents in LaTeX documents
\usepackage{pdfpages}

% defines new commands \Centering, \RaggedLeft, and \RaggedRight and new environments Center, FlushLeft, and FlushRight, which set ragged text and are easily configurable to allow hyphenation
\usepackage{ragged2e}

% makes the page margins visible and hence
% \usepackage{showframe}

% A comprehensive (SI) units package
\usepackage{siunitx}

% provides hyphenatable spacing out (letterspacing), underlining, striking out, etc., using the TeX hyphenation algorithm to find the proper hyphens automatically
  \usepackage{soul}
  \makeatletter
  \let\HL\hl
  \renewcommand\hl{%
    \let\set@color\beamerorig@set@color
    \let\reset@color\beamerorig@reset@color
    \HL}
  \makeatother

% permit use of UTF-8 characters in soul
\usepackage{soulutf8}

% inclusão de Subfiguras https://www.ctan.org/pkg/subfig
\usepackage{subfig}

% for creating graphics programmatically
\usepackage{tikz}
  \usetikzlibrary{arrows,automata,shadows}

% select Adobe Times Roman (or equivalent) as default font
% \usepackage{times}

% Using a color stack for transparency with pdfTEX
\usepackage{transparent}

% provides many glyphs like male and female symbols and astronomical symbols, as well as the complete lasy font set and other odds and ends
\usepackage{wasysym}

% font handwrite
\usepackage{wedn}
  % inline font modify
  \newcommand{\setfont}[2]{{\fontfamily{#1}\selectfont #2}}
  % Usage
  %\setfont{wedn}{Texto...}

% adding colors to your text is supported by the color package
\usepackage{xcolor}
% ------------------------------------------------------------------------------



% ------------------------------------------------------------------------------
% ------------------------------------------------------------------------------
% ------------------------------------------------------------------------------
% ------------------------------------------------------------------------------
% ------------------------------------------------------------------------------
% Pacotes ABNTEX

% citações padrão ABNT
\usepackage[alf]{abntex2cite}
% ------------------------------------------------------------------------------



% ------------------------------------------------------------------------------
% ------------------------------------------------------------------------------
% ------------------------------------------------------------------------------
% ------------------------------------------------------------------------------
% ------------------------------------------------------------------------------
% Macros do documento
\def\universidade{UNINASSAU}

\def\idcurso{31}
\def\curso{Curso Superior de Tecnologia em Análise e Desenvolvimento de Sistemas}

\def\iddisciplina{CODDISC}
\def\disciplina{Nome da Disciplina}

\def\professor{João Ferreira da Silva Júnior}
\def\professorfoto{img/sec01/professor.png}
\def\professoremail{joaoferreirape@gmail.com}
\def\professorsite{https://joaoferreirape.wordpress.com}
\def\professorlattes{http://lattes.cnpq.br/8904695743376784}

\def\semana{}
\def\diadaaula{}
\def\mesdaaula{}
\def\anodaaula{}
\def\semestre{2020-1}
% \def\data{\diadaaula~de~\mesdaaula~de~\anodaaula}
\def\data{Semestre: \semestre}

\title[\iddisciplina]{\textbf{\curso}}

\subtitle{\iddisciplina~---~\disciplina}

\subject{\curso~---~\disciplina~---~\semestre}

\author[SILVA JUNIOR, J. F. $\therefore$]{\href{\professorlattes}{\textbf{\professor}}}

\logo{\includegraphics[height=6ex]{img/logo/brasao.png}}

\institute[UNINASSAU]{\includegraphics[width=40ex]{img/logo/logo.png}}

\date{\data}
% ------------------------------------------------------------------------------



% ------------------------------------------------------------------------------
% ------------------------------------------------------------------------------
% ------------------------------------------------------------------------------
% ------------------------------------------------------------------------------
% ------------------------------------------------------------------------------
% configurações do pacote hyperref, precisa estar aqui devido as macros de dados do documento
\makeatletter
\hypersetup{% http://www.tug.org/applications/hyperref/manual.html#x1-120003.8
  anchorcolor={black},                      % set color of anchors
  % backref=false,                            % do bibliographical back references
  baseurl={\professorlattes},               % set base URL for document
  urlcolor={black},%{blue},                 % url color link
  % bookmarks={true},                         % make bookmarks
  bookmarksnumbered={false},                % put section numbers in bookmarks
  bookmarksopen={true},                     % open up bookmark tree
  bookmarksopenlevel={\maxdimen},           % level to which bookmarks are open
  % bookmarksdepth,                           % <no-value>, <number>, <part,chapter,section,...>
  bookmarkstype={toc},                      % to specify which ‘toc’ file to mimic
  % breaklinks={false},                       % allow links to break over lines
  % citebordercolor=0 1 0,                    % color of border around cites
  citecolor={black},%{green},               % color of citation links
  colorlinks={false},
  % false,                                    % color links
  % true,                                     % (tex4ht, dviwindo)
  % debug=false,                              % provide details of anchors defined; same as verbose
  % destlabel=false,                          % destinations are named by the first \label after the anchor creation
  draft={false},                            % do not do any hyperlinking
  % filebordercolor=0 .5 .5,                  % color of border around file links
  filecolor={black},%{cyan},                % color of file links
  final={true},                             % opposite of option draft
  % frenchlinks=false,                        % use small caps instead of color for links
  % hyperfigures=false,                       % make figures hyper links
  % hyperfootnotes=true,                      % set up hyperlinked footnotes
  % hyperindex=true,                          % set up hyperlinked indices
  % linkbordercolor=1 0 0,                    % color of border around links
  linkcolor={black},%{red},                 % color of links
  % linktocpage={false},                      % make page number, not text, be link on TOC, LOF and LOT
  % menubordercolor=1 0 0,                    % color of border around menu links
  menucolor={black},%{red},                 % color for menu links
  pageanchor={true},                        % put an anchor on every page
  % pagebackref={true},                       % backreference by page number
  pdfauthor={\professor},                   % text for PDF Author field
  % pdfborder=0 0 1,                          % width of PDF link border
  % 0 0 0,                                    % (colorlinks)
  % pdfcenterwindow=true,                     % position the document window in the center of the screen
  pdfcreator={LaTeX with hyperref package, feet: \professor},
  % text for PDF Creator field
  pdfdisplaydoctitle=true,                  % display document title instead of file name in title bar
  pdfkeywords={{\universidade};{\curso};{\disciplina};{\professor}},
  % text for PDF Keywords field
  pdflang={pt-BR},                          % PDF language identifier (RFC 3066)
  pdfproducer={\professor},                 % text for PDF Producer field
  % pdfstartpage=1,                           % page at which PDF document opens
  pdfsubject={\curso~--~\disciplina~--~\semestre},  % text for PDF Subject field
  pdftitle={\curso~--~\disciplina~--~\semestre},    % text for PDF Title field
}
\makeatother
% ------------------------------------------------------------------------------



% ------------------------------------------------------------------------------
% ------------------------------------------------------------------------------
% ------------------------------------------------------------------------------
% ------------------------------------------------------------------------------
% ------------------------------------------------------------------------------
\begin{document}



% ------------------------------------------------------------------------------
% ------------------------------------------------------------------------------
% ------------------------------------------------------------------------------
% ------------------------------------------------------------------------------
% ------------------------------------------------------------------------------
% ------------------------------------------------------------------------------
% ------------------------------------------------------------------------------
% ------------------------------------------------------------------------------
% ------------------------------------------------------------------------------
% ------------------------------------------------------------------------------
\section*{Capa Externa}\label{sec:capa-externa}



{\usebackgroundtemplate{\includegraphics[width=\paperwidth,height=\paperheight,keepaspectratio]{img/logo/background-amarelo.png}}
  \begin{frame}[plain]

    ~

  \end{frame}
} % usebackgroundtemplate



% ------------------------------------------------------------------------------
% ------------------------------------------------------------------------------
% ------------------------------------------------------------------------------
% ------------------------------------------------------------------------------
% ------------------------------------------------------------------------------
% ------------------------------------------------------------------------------
% ------------------------------------------------------------------------------
% ------------------------------------------------------------------------------
% ------------------------------------------------------------------------------
% ------------------------------------------------------------------------------
\section*{Capa}\label{sec:capa}



{\usebackgroundtemplate{{\transparent{0.5}\includegraphics[width=\paperwidth,height=\paperheight,keepaspectratio]{img/logo/background-amarelo.png}}}
  \begin{frame}[plain,t]

    \centering{}

    \addtobeamertemplate{block begin}{\pgfsetfillopacity{0.5}}{\pgfsetfillopacity{1}}
    % \setbeamercolor{block title}{use=structure,fg=white,bg=blue}
    \setbeamercolor{block body}{use=structure,fg=black,bg=blue!20!white}
    \begin{block}{}
      \centering{}
      \huge{\textbf{\MakeUppercase{\curso}}}
    \end{block}

    \vfill

    \vspace*{0.1\paperheight}

    \vfill

    \addtobeamertemplate{block begin}{\pgfsetfillopacity{0.5}}{\pgfsetfillopacity{1}}
    % \setbeamercolor{block title}{use=structure,fg=white,bg=blue}
    \setbeamercolor{block body}{use=structure,fg=black,bg=blue!20!white}
    \begin{block}{}
      \centering{}
      \vspace*{2em}
      \huge{\textbf{\MakeUppercase{\disciplina}}}

      \small{\textbf{\MakeUppercase{\professor}}}
      \vspace*{2em}
    \end{block}

  \end{frame}
} % usebackgroundtemplate



% ------------------------------------------------------------------------------
% ------------------------------------------------------------------------------
% ------------------------------------------------------------------------------
% ------------------------------------------------------------------------------
% ------------------------------------------------------------------------------
% ------------------------------------------------------------------------------
% ------------------------------------------------------------------------------
% ------------------------------------------------------------------------------
% ------------------------------------------------------------------------------
% ------------------------------------------------------------------------------
\section*{Agenda}\label{sec:agenda}



\begin{frame}[allowframebreaks,t]\frametitle{Agenda ---}
% \begin{frame}[t]\frametitle{Agenda}

  \tableofcontents{}
  % \tableofcontents[sectionstyle=show/hide/shaded,subsectionstyle=show/hide/shaded]
  % \tableofcontents[sectionstyle=show,subsectionstyle=show]
  % \tableofcontents[hideallsubsections]
  % \tableofcontents[currentsection,currentsubsection]

\end{frame}



% ------------------------------------------------------------------------------
% ------------------------------------------------------------------------------
% ------------------------------------------------------------------------------
% ------------------------------------------------------------------------------
% ------------------------------------------------------------------------------
% ------------------------------------------------------------------------------
% ------------------------------------------------------------------------------
% ------------------------------------------------------------------------------
% ------------------------------------------------------------------------------
% ------------------------------------------------------------------------------
% !TeX encoding = UTF-8
% !TeX spellcheck = pt_BR
% !TeX root = 20201-CODDISC-JOAO-FERREIRA.tex
% ------------------------------------------------------------------------------

% ------------------------------------------------------------------------------
% Template para aulas da UNINASSAU
% Autor: João Ferreira da Silva Júnior
% Versão: 0.1 - 2020-03-27
% ------------------------------------------------------------------------------



% ------------------------------------------------------------------------------
% ------------------------------------------------------------------------------
% ------------------------------------------------------------------------------
% ------------------------------------------------------------------------------
% ------------------------------------------------------------------------------
% ------------------------------------------------------------------------------
% ------------------------------------------------------------------------------
% ------------------------------------------------------------------------------
% ------------------------------------------------------------------------------
% ------------------------------------------------------------------------------
\section*{Professor}\label{sec:professor}



% \begin{frame}[allowframebreaks,t]\frametitle{Professor ---}
% \begin{frame}[t]\frametitle{Professor}
\begin{frame}[t]

  \begin{columns}

    \begin{column}{0.5\linewidth}
      \begin{figure}[htb]
        \centering{}
        \includegraphics[width=0.5\linewidth,height=0.3\textheight,keepaspectratio]{\professorfoto}

        \textbf{\href{\professorlattes}{\professor}}
      \end{figure}

      \begin{itemize}
        \justifying{}
        \setlength\itemsep{1em}
        \item \href{\professorsite}{\professorsite}
        \item \href{\professorlattes}{\professorlattes}
        \item Formado em Análise e Desenvolvimento de Sistemas em 2013 (UNOPAR).
      \end{itemize}
    \end{column}

    \begin{column}{0.5\linewidth}
      \begin{itemize}
      \justifying{}
      \setlength\itemsep{1em}
      \item Mestre em Engenharia Elétrica, título obtido em 2017 (UFPE).
      \item Doutorando em Engenharia Elétrica, previsão para 07/2021 (UFPE).
      \item Servidor Público na Compesa atuando como Analista de Tecnologia da Informação (Analista de Negócios, Administrador de Banco de Dados, Projetos de Inovação).
      \item Professor da UNINASSAU nos cursos de Análise e Desenvolvimento de Sistemas, Ciência da Computação, e Redes de Computadores.
      \end{itemize}
    \end{column}

  \end{columns}

\end{frame}


% !TeX encoding = UTF-8
% !TeX spellcheck = pt_BR
% !TeX root = 20201-CODDISC-JOAO-FERREIRA.tex
% ------------------------------------------------------------------------------

% ------------------------------------------------------------------------------
% Template para aulas da UNINASSAU
% Autor: João Ferreira da Silva Júnior
% Versão: 0.1 - 2020-03-27
% ------------------------------------------------------------------------------



% ------------------------------------------------------------------------------
% ------------------------------------------------------------------------------
% ------------------------------------------------------------------------------
% ------------------------------------------------------------------------------
% ------------------------------------------------------------------------------
% ------------------------------------------------------------------------------
% ------------------------------------------------------------------------------
% ------------------------------------------------------------------------------
% ------------------------------------------------------------------------------
% ------------------------------------------------------------------------------
\section[Ementa da Disciplina]{Ementa da Disciplina}\label{sec:ementa-disciplina}


{\usebackgroundtemplate{{\transparent{0.2}\includegraphics[width=\paperwidth,height=\paperheight,keepaspectratio]{img/logo/background-branco.png}}}
% \begin{frame}[allowframebreaks,t]%\frametitle{ ---}
  \begin{frame}[plain]%\frametitle{}

    \vfill
    \centering

    \addtobeamertemplate{block begin}{\pgfsetfillopacity{0.5}}{\pgfsetfillopacity{1}}
    % \setbeamercolor{block title}{use=structure,fg=white,bg=blue}
    \setbeamercolor{block body}{use=structure,fg=black,bg=blue!50!white}
    \begin{block}{}
      \centering{}
      \Huge{\textbf{Ementa da Disciplina}}
    \end{block}

    \vfill

\end{frame}
} % usebackgroundtemplate



% ------------------------------------------------------------------------------
% ------------------------------------------------------------------------------
% ------------------------------------------------------------------------------
% ------------------------------------------------------------------------------
% ------------------------------------------------------------------------------
\subsection[Objetivos]{Objetivos}\label{subsec:ementa-objetivos}



% \begin{frame}[allowframebreaks,t]\frametitle{Objetivos ---}
\begin{frame}[t]\frametitle{Objetivos}

  \begin{itemize}
    \justifying{}
    \setlength\itemsep{1em}
    \item Exemplos de citações.
    \item \cite{Tanenbaum2011}
    \item \cite{Tanenbaum2016}
  \end{itemize}

\end{frame}



% ------------------------------------------------------------------------------
% ------------------------------------------------------------------------------
% ------------------------------------------------------------------------------
% ------------------------------------------------------------------------------
% ------------------------------------------------------------------------------
\subsection[Ementa]{Ementa}\label{subsec:ementa-ementa}



% \begin{frame}[allowframebreaks,t]\frametitle{Ementa ---}
\begin{frame}[t]\frametitle{Ementa}

  \begin{itemize}
    \justifying{}
    \setlength\itemsep{1em}
    \item Item 1.
    \item Item 2.
    \item Item 3.
  \end{itemize}

\end{frame}



% ------------------------------------------------------------------------------
% ------------------------------------------------------------------------------
% ------------------------------------------------------------------------------
% ------------------------------------------------------------------------------
% ------------------------------------------------------------------------------
\subsection[Metodologia de Ensino]{Metodologia de Ensino}\label{subsec:ementa-metodologia-ensino}



% \begin{frame}[allowframebreaks,t]\frametitle{Metodologia de Ensino ---}
\begin{frame}[t]\frametitle{Metodologia de Ensino}

  \begin{itemize}
    \justifying{}
    \setlength\itemsep{1em}
    \item Item 1.
    \item Item 2.
    \item Item 3.
  \end{itemize}

\end{frame}



% ------------------------------------------------------------------------------
% ------------------------------------------------------------------------------
% ------------------------------------------------------------------------------
% ------------------------------------------------------------------------------
% ------------------------------------------------------------------------------
\subsection[Metodologia de Avaliação]{Metodologia de Avaliação}\label{subsec:ementa-metodologia-avaliacao}



% \begin{frame}[allowframebreaks,t]\frametitle{Metodologia de Avaliação ---}
\begin{frame}[t]\frametitle{Metodologia de Avaliação}

  \begin{itemize}
    \justifying{}
    \setlength\itemsep{1em}
    \item Item 1.
    \item Item 2.
    \item Item 3.
  \end{itemize}

\end{frame}



% ------------------------------------------------------------------------------
% ------------------------------------------------------------------------------
% ------------------------------------------------------------------------------
% ------------------------------------------------------------------------------
% ------------------------------------------------------------------------------
\subsection[Bibliografia]{Bibliografia}\label{subsec:ementa-bibliografia}



% \begin{frame}[allowframebreaks,t]\frametitle{Bibliografia Básica ---}
\begin{frame}[t]\frametitle{Bibliografia Básica}

  \begin{itemize}
    \justifying{}
    \setlength\itemsep{1em}
    \item Item 1.
    \item Item 2.
    \item Item 3.
  \end{itemize}

\end{frame}



% \begin{frame}[allowframebreaks,t]\frametitle{Bibliografia Complementar ---}
\begin{frame}[t]\frametitle{Bibliografia Complementar}

  \begin{itemize}
    \justifying{}
    \setlength\itemsep{1em}
    \item Item 1.
    \item Item 2.
    \item Item 3.
  \end{itemize}

\end{frame}


% !TeX encoding = UTF-8
% !TeX spellcheck = pt_BR
% !TeX root = 20201-CODDISC-JOAO-FERREIRA.tex
% ------------------------------------------------------------------------------

% ------------------------------------------------------------------------------
% Template para aulas da UNINASSAU
% Autor: João Ferreira da Silva Júnior
% Versão: 0.1 - 2020-03-27
% ------------------------------------------------------------------------------



% ------------------------------------------------------------------------------
% ------------------------------------------------------------------------------
% ------------------------------------------------------------------------------
% ------------------------------------------------------------------------------
% ------------------------------------------------------------------------------
% ------------------------------------------------------------------------------
% ------------------------------------------------------------------------------
% ------------------------------------------------------------------------------
% ------------------------------------------------------------------------------
% ------------------------------------------------------------------------------
\section[Planejamento da Disciplina]{Planejamento da Disciplina}\label{sec:planejamento}



{\usebackgroundtemplate{{\transparent{0.5}\includegraphics[width=\paperwidth,height=\paperheight,keepaspectratio]{img/logo/background-branco.png}}}
  % \begin{frame}[allowframebreaks,t]%\frametitle{ ---}
  \begin{frame}[plain]%\frametitle{}

    \vfill
    \centering

    \addtobeamertemplate{block begin}{\pgfsetfillopacity{0.5}}{\pgfsetfillopacity{1}}
    % \setbeamercolor{block title}{use=structure,fg=white,bg=blue}
    \setbeamercolor{block body}{use=structure,fg=black,bg=blue!20!white}
    \begin{block}{}
      \centering{}
      \Huge{\textbf{Planejamento da Disciplina}}
    \end{block}

    \vfill

  \end{frame}
} % usebackgroundtemplate



% ------------------------------------------------------------------------------
% ------------------------------------------------------------------------------
% ------------------------------------------------------------------------------
% ------------------------------------------------------------------------------
% ------------------------------------------------------------------------------
\subsection[Calendário Acadêmico]{Calendário Acadêmico}\label{subsec:planejamento-calendario-academico}



% \begin{frame}[allowframebreaks,t]\frametitle{Calendário Acadêmico ---}
\begin{frame}[t]\frametitle{Calendário Acadêmico}

  \begin{itemize}
    \justifying{}
    \setlength\itemsep{1em}
    \item Objetivo 1.
    \item Objetivo 2.
    \item Objetivo 3.
  \end{itemize}

\end{frame}



% ------------------------------------------------------------------------------
% ------------------------------------------------------------------------------
% ------------------------------------------------------------------------------
% ------------------------------------------------------------------------------
% ------------------------------------------------------------------------------
\subsection[Participação]{Participação}\label{subsec:planejamento-participacao}



% \begin{frame}[allowframebreaks,t]\frametitle{Participação ---}
\begin{frame}[t]\frametitle{Participação}

  \begin{itemize}
    \justifying{}
    \setlength\itemsep{1em}
    \item Item 1.
    \item Item 2.
    \item Item 3.
  \end{itemize}

\end{frame}



% ------------------------------------------------------------------------------
% ------------------------------------------------------------------------------
% ------------------------------------------------------------------------------
% ------------------------------------------------------------------------------
% ------------------------------------------------------------------------------
\subsection[Lista de Exercícios]{Lista de Exercícios}\label{subsec:planejamento-lista-exercicios}



% \begin{frame}[allowframebreaks,t]\frametitle{Lista de Exercícios ---}
\begin{frame}[t]\frametitle{Lista de Exercícios}

  \begin{itemize}
    \justifying{}
    \setlength\itemsep{1em}
    \item Item 1.
    \item Item 2.
    \item Item 3.
  \end{itemize}

\end{frame}



% ------------------------------------------------------------------------------
% ------------------------------------------------------------------------------
% ------------------------------------------------------------------------------
% ------------------------------------------------------------------------------
% ------------------------------------------------------------------------------
\subsection[Avaliações]{Avaliações}\label{subsec:planejamento-avaliacoes}



% \begin{frame}[allowframebreaks,t]\frametitle{Avaliações ---}
\begin{frame}[t]\frametitle{Avaliações}

  \begin{itemize}
    \justifying{}
    \setlength\itemsep{1em}
    \item Item 1.
    \item Item 2.
    \item Item 3.
  \end{itemize}

\end{frame}



% ------------------------------------------------------------------------------
% ------------------------------------------------------------------------------
% ------------------------------------------------------------------------------
% ------------------------------------------------------------------------------
% ------------------------------------------------------------------------------
\subsection[Comunicação da Turma]{Comunicação da Turma}\label{subsec:planejamento-comunicacao-turma}



% \begin{frame}[allowframebreaks,t]\frametitle{Comunicação da Turma ---}
\begin{frame}[t]\frametitle{Comunicação da Turma}

  \begin{figure}[htb]
    \centering{}
    \includegraphics[width=0.07\paperwidth]{img/sec03/telegram.png}

    \textbf{\href{https://telegram.org}{Telegram}}
  \end{figure}

  \begin{columns}
    \begin{column}{0.45\linewidth}
      \textbf{Vantagens:}
      \begin{itemize}
        \justifying{}
        \setlength\itemsep{1em}
        \item Grupo pode ter muitos membros e a troca de arquivos é bem otimizada.
        \item Histórico permanente, quem entre depois poderá ver todas as mensagens já trocadas no grupo.
        \item Ótimo aplicativo para desktop/web.
      \end{itemize}
    \end{column}

    \begin{column}{0.45\linewidth}
      \textbf{Disponível para:}
      \begin{itemize}
        \justifying{}
        % \setlength\itemsep{1em}
        \item \href{https://telegram.org/dl/android}{Android}
        \item \href{https://telegram.org/dl/ios}{iPhone/iPad}
        \item \href{https://desktop.telegram.org}{PC/Mac/Linux}
        \item \href{https://web.telegram.org}{Web}
        \item \href{https://macos.telegram.org}{macOS}
      \end{itemize}

      \textbf{Link para participar do grupo:}

      \href{https://web.telegram.org}{https://web.telegram.org}
    \end{column}
  \end{columns}

\end{frame}


% !TeX encoding = UTF-8
% !TeX spellcheck = pt_BR
% !TeX root = 20201-CODDISC-JOAO-FERREIRA.tex
% ------------------------------------------------------------------------------

% ------------------------------------------------------------------------------
% Template para aulas da UNINASSAU
% Autor: João Ferreira da Silva Júnior
% Versão: 0.1 - 2020-03-27
% ------------------------------------------------------------------------------



% ------------------------------------------------------------------------------
% ------------------------------------------------------------------------------
% ------------------------------------------------------------------------------
% ------------------------------------------------------------------------------
% ------------------------------------------------------------------------------
% ------------------------------------------------------------------------------
% ------------------------------------------------------------------------------
% ------------------------------------------------------------------------------
% ------------------------------------------------------------------------------
% ------------------------------------------------------------------------------
\section[Introdução da Disciplina]{Introdução da Disciplina}\label{sec:introducao}



{\usebackgroundtemplate{{\transparent{0.2}\includegraphics[width=\paperwidth,height=\paperheight,keepaspectratio]{img/logo/background-branco.png}}}
  % \begin{frame}[allowframebreaks,t]%\frametitle{ ---}
  \begin{frame}[plain]%\frametitle{}

    \vfill
    \centering

    \addtobeamertemplate{block begin}{\pgfsetfillopacity{0.5}}{\pgfsetfillopacity{1}}
    % \setbeamercolor{block title}{use=structure,fg=white,bg=blue}
    \setbeamercolor{block body}{use=structure,fg=black,bg=blue!50!white}
    \begin{block}{}
      \centering{}
      \Huge{\textbf{Introdução da Disciplina}}
    \end{block}

    \vfill

  \end{frame}
} % usebackgroundtemplate



% ------------------------------------------------------------------------------
% ------------------------------------------------------------------------------
% ------------------------------------------------------------------------------
% ------------------------------------------------------------------------------
% ------------------------------------------------------------------------------
\subsection[Tópico 1]{Tópico 1}\label{subsec:introducao-topico1}



% \begin{frame}[allowframebreaks,t]\frametitle{Tópico 1 ---}
\begin{frame}[t]\frametitle{Tópico 1}

  \begin{itemize}
    \justifying{}
    \setlength\itemsep{1em}
    \item Item 1.
    \item Item 2.
    \item Item 3.
  \end{itemize}

\end{frame}



% ------------------------------------------------------------------------------
% ------------------------------------------------------------------------------
% ------------------------------------------------------------------------------
% ------------------------------------------------------------------------------
% ------------------------------------------------------------------------------
\subsection[Exercício 1]{Exercício 1}\label{subsec:exercicio1}



% \begin{frame}[allowframebreaks,t]\frametitle{Tópico 1 ---}
\begin{frame}[t]\frametitle{Tópico 1}

\begin{itemize}
  \justifying{}
  \setlength\itemsep{1em}
  \item Questão 1.
  \item Questão 2.
  \item Questão 3.
\end{itemize}

\end{frame}


% !TeX encoding = UTF-8
% !TeX spellcheck = pt_BR
% !TeX root = 20201-CODDISC-JOAO-FERREIRA.tex
% ------------------------------------------------------------------------------

% ------------------------------------------------------------------------------
% Template para aulas da UNINASSAU
% Autor: João Ferreira da Silva Júnior
% Versão: 0.1 - 2020-03-27
% ------------------------------------------------------------------------------



% ------------------------------------------------------------------------------
% ------------------------------------------------------------------------------
% ------------------------------------------------------------------------------
% ------------------------------------------------------------------------------
% ------------------------------------------------------------------------------
% ------------------------------------------------------------------------------
% ------------------------------------------------------------------------------
% ------------------------------------------------------------------------------
% ------------------------------------------------------------------------------
% ------------------------------------------------------------------------------
\section[Tema 2 da Disciplina]{Tema 2 da Disciplina}\label{sec:tema-2}



{\usebackgroundtemplate{{\transparent{0.2}\includegraphics[width=\paperwidth,height=\paperheight,keepaspectratio]{img/logo/background-branco.png}}}
  % \begin{frame}[allowframebreaks,t]%\frametitle{ ---}
  \begin{frame}[plain]%\frametitle{}

    \vfill
    \centering

    \addtobeamertemplate{block begin}{\pgfsetfillopacity{0.5}}{\pgfsetfillopacity{1}}
    % \setbeamercolor{block title}{use=structure,fg=white,bg=blue}
    \setbeamercolor{block body}{use=structure,fg=black,bg=blue!50!white}
    \begin{block}{}
      \centering{}
      \Huge{\textbf{Tema 2 da Disciplina}}
    \end{block}

    \vfill

  \end{frame}
} % usebackgroundtemplate



% ------------------------------------------------------------------------------
% ------------------------------------------------------------------------------
% ------------------------------------------------------------------------------
% ------------------------------------------------------------------------------
% ------------------------------------------------------------------------------
\subsection[Tópico 1]{Tópico 1}\label{subsec:tema-2-topico-1}



% \begin{frame}[allowframebreaks,t]\frametitle{Tópico 1 ---}
\begin{frame}[t]\frametitle{Tópico 1}

  \begin{itemize}
    \justifying{}
    \setlength\itemsep{1em}
    \item Item 1.
    \item Item 2.
    \item Item 3.
  \end{itemize}

\end{frame}



% ------------------------------------------------------------------------------
% ------------------------------------------------------------------------------
% ------------------------------------------------------------------------------
% ------------------------------------------------------------------------------
% ------------------------------------------------------------------------------
\subsection[Exercício 2]{Exercício 2}\label{subsec:exercicio-2}



% \begin{frame}[allowframebreaks,t]\frametitle{Tópico 1 ---}
\begin{frame}[t]\frametitle{Tópico 1}

\begin{itemize}
  \justifying{}
  \setlength\itemsep{1em}
  \item Questão 1.
  \item Questão 2.
  \item Questão 3.
\end{itemize}

\end{frame}


% !TeX encoding = UTF-8
% !TeX spellcheck = pt_BR
% !TeX root = 20201-CODDISC-JOAO-FERREIRA.tex
% ------------------------------------------------------------------------------

% ------------------------------------------------------------------------------
% Template para aulas da UNINASSAU
% Autor: João Ferreira da Silva Júnior
% Versão: 0.1 - 2020-03-27
% ------------------------------------------------------------------------------



% ------------------------------------------------------------------------------
% ------------------------------------------------------------------------------
% ------------------------------------------------------------------------------
% ------------------------------------------------------------------------------
% ------------------------------------------------------------------------------
% ------------------------------------------------------------------------------
% ------------------------------------------------------------------------------
% ------------------------------------------------------------------------------
% ------------------------------------------------------------------------------
% ------------------------------------------------------------------------------
\section[Lista de Exercícios 1]{Lista de Exercícios 1}\label{sec:lista-exercicios-1}



{\usebackgroundtemplate{{\transparent{0.5}\includegraphics[width=\paperwidth,height=\paperheight,keepaspectratio]{img/logo/background-branco.png}}}
  % \begin{frame}[allowframebreaks,t]%\frametitle{ ---}
  \begin{frame}[plain]%\frametitle{}

    \vfill
    \centering

    \addtobeamertemplate{block begin}{\pgfsetfillopacity{0.5}}{\pgfsetfillopacity{1}}
    % \setbeamercolor{block title}{use=structure,fg=white,bg=blue}
    \setbeamercolor{block body}{use=structure,fg=black,bg=blue!20!white}
    \begin{block}{}
      \centering{}
      \Huge{\textbf{Lista de Exercícios 1}}
    \end{block}

    \vfill

  \end{frame}
} % usebackgroundtemplate



% ------------------------------------------------------------------------------
% ------------------------------------------------------------------------------
% ------------------------------------------------------------------------------
% ------------------------------------------------------------------------------
% ------------------------------------------------------------------------------
\subsection[Exercício 1]{Exercício 1}\label{subsec:lista-exercicios-1}



% \begin{frame}[allowframebreaks,t]\frametitle{Tópico 1 ---}
\begin{frame}[t]\frametitle{Tópico 1}

\begin{itemize}
  \justifying{}
  \setlength\itemsep{1em}
  \item Questão 1.
  \item Questão 2.
  \item Questão 3.
\end{itemize}

\end{frame}


% !TeX encoding = UTF-8
% !TeX spellcheck = pt_BR
% !TeX root = 20201-CODDISC-JOAO-FERREIRA.tex
% ------------------------------------------------------------------------------

% ------------------------------------------------------------------------------
% Template para aulas da UNINASSAU
% Autor: João Ferreira da Silva Júnior
% Versão: 0.1 - 2020-03-27
% ------------------------------------------------------------------------------



% ------------------------------------------------------------------------------
% ------------------------------------------------------------------------------
% ------------------------------------------------------------------------------
% ------------------------------------------------------------------------------
% ------------------------------------------------------------------------------
% ------------------------------------------------------------------------------
% ------------------------------------------------------------------------------
% ------------------------------------------------------------------------------
% ------------------------------------------------------------------------------
% ------------------------------------------------------------------------------
\section[Tema 3 da Disciplina]{Tema 3 da Disciplina}\label{sec:tema-3}



{\usebackgroundtemplate{{\transparent{0.5}\includegraphics[width=\paperwidth,height=\paperheight,keepaspectratio]{img/logo/background-branco.png}}}
  % \begin{frame}[allowframebreaks,t]%\frametitle{ ---}
  \begin{frame}[plain]%\frametitle{}

    \vfill
    \centering

    \addtobeamertemplate{block begin}{\pgfsetfillopacity{0.5}}{\pgfsetfillopacity{1}}
    % \setbeamercolor{block title}{use=structure,fg=white,bg=blue}
    \setbeamercolor{block body}{use=structure,fg=black,bg=blue!20!white}
    \begin{block}{}
      \centering{}
      \Huge{\textbf{Tema 3 da Disciplina}}
    \end{block}

    \vfill

  \end{frame}
} % usebackgroundtemplate



% ------------------------------------------------------------------------------
% ------------------------------------------------------------------------------
% ------------------------------------------------------------------------------
% ------------------------------------------------------------------------------
% ------------------------------------------------------------------------------
\subsection[Tópico 1]{Tópico 1}\label{subsec:tema-3-topico1}



% \begin{frame}[allowframebreaks,t]\frametitle{Tópico 1 ---}
\begin{frame}[t]\frametitle{Tópico 1}

  \begin{itemize}
    \justifying{}
    \setlength\itemsep{1em}
    \item Item 1.
    \item Item 2.
    \item Item 3.
  \end{itemize}

\end{frame}



% ------------------------------------------------------------------------------
% ------------------------------------------------------------------------------
% ------------------------------------------------------------------------------
% ------------------------------------------------------------------------------
% ------------------------------------------------------------------------------
\subsection[Exercício 3]{Exercício 3}\label{subsec:exercicio-3}



% \begin{frame}[allowframebreaks,t]\frametitle{Tópico 1 ---}
\begin{frame}[t]\frametitle{Tópico 1}

\begin{itemize}
  \justifying{}
  \setlength\itemsep{1em}
  \item Questão 1.
  \item Questão 2.
  \item Questão 3.
\end{itemize}

\end{frame}


% !TeX encoding = UTF-8
% !TeX spellcheck = pt_BR
% !TeX root = 20201-CODDISC-JOAO-FERREIRA.tex
% ------------------------------------------------------------------------------

% ------------------------------------------------------------------------------
% Template para aulas da UNINASSAU
% Autor: João Ferreira da Silva Júnior
% Versão: 0.1 - 2020-03-27
% ------------------------------------------------------------------------------



% ------------------------------------------------------------------------------
% ------------------------------------------------------------------------------
% ------------------------------------------------------------------------------
% ------------------------------------------------------------------------------
% ------------------------------------------------------------------------------
% ------------------------------------------------------------------------------
% ------------------------------------------------------------------------------
% ------------------------------------------------------------------------------
% ------------------------------------------------------------------------------
% ------------------------------------------------------------------------------
\section[Tema 4 da Disciplina]{Tema 4 da Disciplina}\label{sec:tema-4}



{\usebackgroundtemplate{{\transparent{0.2}\includegraphics[width=\paperwidth,height=\paperheight,keepaspectratio]{img/logo/background-branco.png}}}
  % \begin{frame}[allowframebreaks,t]%\frametitle{ ---}
  \begin{frame}[plain]%\frametitle{}

    \vfill
    \centering

    \addtobeamertemplate{block begin}{\pgfsetfillopacity{0.5}}{\pgfsetfillopacity{1}}
    % \setbeamercolor{block title}{use=structure,fg=white,bg=blue}
    \setbeamercolor{block body}{use=structure,fg=black,bg=blue!50!white}
    \begin{block}{}
      \centering{}
      \Huge{\textbf{Tema 4 da Disciplina}}
    \end{block}

    \vfill

  \end{frame}
} % usebackgroundtemplate



% ------------------------------------------------------------------------------
% ------------------------------------------------------------------------------
% ------------------------------------------------------------------------------
% ------------------------------------------------------------------------------
% ------------------------------------------------------------------------------
\subsection[Tópico 1]{Tópico 1}\label{subsec:tema-4-topico1}



% \begin{frame}[allowframebreaks,t]\frametitle{Tópico 1 ---}
\begin{frame}[t]\frametitle{Tópico 1}

  \begin{itemize}
    \justifying{}
    \setlength\itemsep{1em}
    \item Item 1.
    \item Item 2.
    \item Item 3.
  \end{itemize}

\end{frame}



% ------------------------------------------------------------------------------
% ------------------------------------------------------------------------------
% ------------------------------------------------------------------------------
% ------------------------------------------------------------------------------
% ------------------------------------------------------------------------------
\subsection[Exercício 4]{Exercício 4}\label{subsec:exercicio-4}



% \begin{frame}[allowframebreaks,t]\frametitle{Tópico 1 ---}
\begin{frame}[t]\frametitle{Tópico 1}

\begin{itemize}
  \justifying{}
  \setlength\itemsep{1em}
  \item Questão 1.
  \item Questão 2.
  \item Questão 3.
\end{itemize}

\end{frame}


% !TeX encoding = UTF-8
% !TeX spellcheck = pt_BR
% !TeX root = 20201-CODDISC-JOAO-FERREIRA.tex
% ------------------------------------------------------------------------------

% ------------------------------------------------------------------------------
% Template para aulas da UNINASSAU
% Autor: João Ferreira da Silva Júnior
% Versão: 0.1 - 2020-03-27
% ------------------------------------------------------------------------------



% ------------------------------------------------------------------------------
% ------------------------------------------------------------------------------
% ------------------------------------------------------------------------------
% ------------------------------------------------------------------------------
% ------------------------------------------------------------------------------
% ------------------------------------------------------------------------------
% ------------------------------------------------------------------------------
% ------------------------------------------------------------------------------
% ------------------------------------------------------------------------------
% ------------------------------------------------------------------------------
\section[Lista de Exercícios 2]{Lista de Exercícios 2}\label{sec:lista-exercicios-2}



{\usebackgroundtemplate{{\transparent{0.5}\includegraphics[width=\paperwidth,height=\paperheight,keepaspectratio]{img/logo/background-branco.png}}}
  % \begin{frame}[allowframebreaks,t]%\frametitle{ ---}
  \begin{frame}[plain]%\frametitle{}

    \vfill
    \centering

    \addtobeamertemplate{block begin}{\pgfsetfillopacity{0.5}}{\pgfsetfillopacity{1}}
    % \setbeamercolor{block title}{use=structure,fg=white,bg=blue}
    \setbeamercolor{block body}{use=structure,fg=black,bg=blue!20!white}
    \begin{block}{}
      \centering{}
      \Huge{\textbf{Lista de Exercícios 2}}
    \end{block}

    \vfill

  \end{frame}
} % usebackgroundtemplate



% ------------------------------------------------------------------------------
% ------------------------------------------------------------------------------
% ------------------------------------------------------------------------------
% ------------------------------------------------------------------------------
% ------------------------------------------------------------------------------
\subsection[Exercício 2]{Exercício 2}\label{subsec:lista-exercicios-2}



% \begin{frame}[allowframebreaks,t]\frametitle{Tópico 1 ---}
\begin{frame}[t]\frametitle{Tópico 1}

\begin{itemize}
  \justifying{}
  \setlength\itemsep{1em}
  \item Questão 1.
  \item Questão 2.
  \item Questão 3.
\end{itemize}

\end{frame}




% ------------------------------------------------------------------------------
% ------------------------------------------------------------------------------
% ------------------------------------------------------------------------------
% ------------------------------------------------------------------------------
% ------------------------------------------------------------------------------
% ------------------------------------------------------------------------------
% ------------------------------------------------------------------------------
% ------------------------------------------------------------------------------
% ------------------------------------------------------------------------------
% ------------------------------------------------------------------------------
\section[Agradecimento]{Agradecimento}\label{sec:agradecimentos}



{\usebackgroundtemplate{{\transparent{0.5}\includegraphics[width=\paperwidth,height=\paperheight,keepaspectratio]{img/logo/background-amarelo.png}}}
  \begin{frame}[plain,t]

    \centering

    \addtobeamertemplate{block begin}{\pgfsetfillopacity{0.5}}{\pgfsetfillopacity{1}}
    % \setbeamercolor{block title}{use=structure,fg=white,bg=blue}
    \setbeamercolor{block body}{use=structure,fg=black,bg=blue!20!white}
    \begin{block}{}
      \centering{}
      \Huge{\textbf{Obrigado!}}
    \end{block}

    \vfill

    \vspace*{0.3\paperheight}

    \vfill

    \addtobeamertemplate{block begin}{\pgfsetfillopacity{0.5}}{\pgfsetfillopacity{1}}
    \setbeamercolor{block title}{use=structure,fg=white,bg=blue}
    \setbeamercolor{block body}{use=structure,fg=black,bg=blue!20!white}
    \begin{block}{\textbf{\professor}}
      \justifying{}
      \begin{itemize}
        \justifying{}
        \setlength\itemsep{1em}
        \item \href{mailto:\professoremail}{\professoremail}
        \item \href{\professorlattes}{\professorlattes}
        \item \href{professorsite}{\professorsite}
      \end{itemize}
    \end{block}

  \end{frame}
} % usebackgroundtemplate



% ------------------------------------------------------------------------------
% ------------------------------------------------------------------------------
% ------------------------------------------------------------------------------
% ------------------------------------------------------------------------------
% ------------------------------------------------------------------------------
% ------------------------------------------------------------------------------
% ------------------------------------------------------------------------------
% ------------------------------------------------------------------------------
% ------------------------------------------------------------------------------
% ------------------------------------------------------------------------------
\section[Referências]{Referências}\label{sec:referencias}



{\usebackgroundtemplate{{\transparent{0.5}\includegraphics[width=\paperwidth,height=\paperheight,keepaspectratio]{img/logo/background-branco.png}}}
  \begin{frame}[plain,allowframebreaks,t]%\frametitle{Referências ---}
  % \begin{frame}[plain,t]%\frametitle{Referências}
    % \bibliographystyle{apalike}
    \bibliography{library.bib}
  \end{frame}
} % usebackgroundtemplate



% ------------------------------------------------------------------------------
% ------------------------------------------------------------------------------
% ------------------------------------------------------------------------------
% ------------------------------------------------------------------------------
% ------------------------------------------------------------------------------
% ------------------------------------------------------------------------------
% ------------------------------------------------------------------------------
% ------------------------------------------------------------------------------
% ------------------------------------------------------------------------------
% ------------------------------------------------------------------------------
\section*{Capa Final}\label{sec:capa-final}



{\usebackgroundtemplate{\includegraphics[width=\paperwidth,height=\paperheight,keepaspectratio]{img/logo/background-branco.png}}
  \begin{frame}[plain]

    ~

  \end{frame}
} % usebackgroundtemplate



\end{document}
